% Responsable: Fernando Rueda
\section{Bloques 4 y 5: modelos, pronóstico y recomendaciones}

Las métricas de esta sección son las que reportamos en los dos cuadernos de modelado del
Laboratorio 1, el de regiones y el de fronteras. Los modelos se
entrenaron con enero de 2009 a marzo de 2021 y se evaluaron con abril de 2021 a junio de 2026,
63 meses. El bloque 4 muestra la serie total, que es la que resume el país y la que un directivo
necesita anticipar.

\subsection{Bloque 4, izquierda: duelo de modelos}

\paragraph{Qué hay que comunicar.} En el laboratorio comparamos seis algoritmos, suavizamiento
exponencial simple, Holt-Winters, Holt-Winters amortiguado, seasonal naive, SARIMA y Prophet, con
dos formas de evaluar. En la primera cada modelo pronostica los 63 meses de prueba de una sola
vez, y todos fallan, porque aprendieron un entrenamiento que termina en el piso de la pandemia y
proyectan esa caída hacia adelante. En la segunda el modelo se actualiza cada mes con el dato
real y pronostica solo el mes siguiente, que es como se usaría en la práctica. El hallazgo para
el directivo no es cuál algoritmo gana por poco, sino que un pronóstico solo sirve si se
actualiza.

\begin{center}
\begin{tabular}{@{}llrrr@{}}
\toprule
Evaluación & Modelo & MAE & RMSE & AIC \\
\midrule
63 meses de una vez & Suavizamiento exponencial simple & 158\,777 & 173\,709 & -- \\
                    & Holt-Winters amortiguado         & 170\,053 & 183\,865 & -- \\
                    & Holt-Winters                     & 178\,392 & 192\,825 & -- \\
                    & Seasonal naive                   & 207\,322 & 219\,664 & -- \\
                    & SARIMA                           & 226\,894 & 243\,234 & 64.1 \\
                    & Prophet                          & 232\,826 & 246\,743 & -- \\
\midrule
Un mes adelante     & \textbf{SARIMA (ganador)}        & \textbf{50\,430} & \textbf{80\,789} & 64.1 \\
                    & Seasonal naive                   & 65\,996 & 93\,379 & -- \\
\bottomrule
\end{tabular}

{\footnotesize Serie total, error en visitantes por mes. El SARIMA es el
SARIMA(0,1,1)(1,1,0)$_{12}$ sobre el logaritmo.}
\end{center}

\paragraph{Qué métrica se muestra.} El gráfico usa el MAE como medida principal porque está en
las mismas unidades que la serie y se traduce sin esfuerzo, ``el modelo se equivoca en promedio
unos 50 mil visitantes por mes''. El RMSE va como cifra secundaria junto a cada barra, porque
castiga más los errores grandes y muestra que el ganador concentra sus fallos en los meses
atípicos de la reapertura. El AIC no se grafica. En el laboratorio sirvió para elegir el mejor
SARIMA entre nueve candidatos, pero no permite comparar modelos de familias distintas ni modelos
ajustados sobre escalas distintas, y el SARIMA trabaja sobre el logaritmo. Ponerlo junto al MAE
invitaría a comparar cifras que no son comparables, así que queda en la ficha técnica y en este
documento.

\paragraph{Tipo de gráfico.} El dato es una métrica numérica para ocho variantes de modelo, una
variable cuantitativa sobre categorías. Usamos barras horizontales ordenadas de menor a mayor
error, porque la longitud desde una base común es la comparación más precisa para el ojo y la
orientación horizontal deja leer el nombre de cada modelo sin girarlo. Las barras parten de cero,
porque recortar el eje haría que una diferencia de 10\% pareciera del doble. Las barras se separan
en dos grupos con una línea fina y un rótulo, ``pronóstico de cinco años de golpe'' y
``actualizado cada mes'', y por color, las del primer grupo en gris claro, la del seasonal naive
actualizado en azul principal y la del ganador en naranja. La diferencia de longitud entre los
dos grupos, más de tres veces, es el mensaje y se ve sin leer las cifras.

\paragraph{Cómo se dice en la página.} El título del gráfico enuncia la conclusión, ``Actualizado
cada mes, el pronóstico erra 3 veces menos''. Debajo, una línea traduce la cifra, el ganador se
equivoca en unos 50 mil visitantes, alrededor de una quinta parte del promedio mensual de la
prueba, que fue de 241 mil.

\subsection{Bloque 4, derecha: pronóstico del modelo ganador}

\paragraph{Tipo de gráfico.} Es una serie de tiempo con incertidumbre, así que usamos un gráfico
de líneas con un abanico. Lo observado va en azul principal y la estimación del SARIMA en naranja,
con dos bandas del mismo tono, la más oscura para el intervalo del 80\% y la más clara para el del
95\%. Las bandas se rotulan con la palabra ``Estimación'', como pide el enunciado, y con una
etiqueta directa que dice qué porcentaje cubre cada una, sin leyenda aparte. Mostrar el intervalo
es tan importante como mostrar la línea, porque le dice al directivo cuánto puede confiar en el
número y evita que lea un pronóstico como una certeza.

\paragraph{Qué periodo.} El gráfico muestra el conjunto de prueba, de abril de 2021 a junio de
2026, con la estimación a un mes adelante, que es exactamente lo que evaluamos en el laboratorio.
Para responder la segunda mitad de la pregunta del enunciado, qué debemos anticipar, se prolonga
con línea punteada la estimación de los doce meses siguientes, de julio de 2026 a junio de 2027,
con el mismo modelo ajustado con toda la serie. Una línea vertical separa lo evaluado de lo
proyectado, y la parte proyectada lleva su propio rótulo para que no se confunda con un resultado
validado.

\paragraph{Decisiones para que el gráfico no mienta.} Como el modelo se ajusta sobre el
logaritmo, la banda es asimétrica, más ancha hacia arriba que hacia abajo, y la dejamos así en
lugar de dibujar una banda simétrica que no corresponde al modelo. El eje vertical empieza en
cero. En los meses de la reapertura de 2022 la banda del 95\% se dispara muy por encima del resto
de la serie, y en lugar de estirar el eje, lo que aplastaría todo el gráfico, se recorta y se
indica con una nota. Así el lector ve la incertidumbre real de esos meses sin perder la lectura
del resto.

\subsection{Bloque 5: recomendaciones y ficha técnica}

\paragraph{Formato.} Las recomendaciones van en cuatro tarjetas del mismo tamaño, como en el
boceto, para que ninguna parezca más importante por su posición y para que se lean de un vistazo.
Cada tarjeta tiene una acción que empieza con un verbo, en azul principal y negrita, y debajo la
cifra que la justifica, con el número en naranja. Una recomendación sin evidencia sería una
opinión, así que la regla es que cada una debe poder rastrearse hasta un gráfico de la página.

\begin{center}
\begin{tabularx}{\textwidth}{@{}clXl@{}}
\toprule
\# & Acción & Evidencia & Bloque \\
\midrule
1 & Reportar a los visitantes extranjeros con una serie reconciliada &
Sin los residentes de Guatemala que dejaron de contarse en 2023, los visitantes de 2025 son el
138\% de los de 2019, mientras que el total crudo marca 81\%. & 1 y 2 \\
2 & Programar campañas y operativos con el calendario de cada mercado &
Centroamérica llena diciembre (1.41), Norteamérica julio (1.32) y Europa noviembre (1.25), cada
una frente a un mes típico. & 3 \\
3 & Reforzar los puestos fronterizos con El Salvador &
San Cristóbal es la serie que más crece, 10.1\% anual en extranjeros entre 2009 y 2025, y la más
volátil, y Valle Nuevo fue la más golpeada por la pandemia ($-$84\% en 2020 y 2021). & 3 \\
4 & Actualizar los pronósticos cada mes &
El mejor pronóstico de una sola vez se equivoca en 159 mil visitantes por mes, el actualizado en
50 mil. & 4 \\
\bottomrule
\end{tabularx}
\end{center}

El orden va de lo más urgente a lo más operativo. La primera corrige cómo se lee toda la
información y por eso va antes que las demás, porque sin ella el directivo creería que el turismo
no se ha recuperado. La segunda y la tercera usan el diagnóstico del bloque 3, cuándo viene la
gente y por dónde. La cuarta cierra con el bloque 4 y es la que permite revisar las otras tres
cada mes.

\paragraph{Ficha técnica.} Va al pie de la página, en el tamaño más pequeño y en gris oscuro,
porque se consulta pero no compite con el mensaje. Incluye la fuente oficial, Base de Migración
2009--2026 de INGUAT e IGM, el periodo de enero de 2009 a junio de 2026, la definición de
visitante como turista más excursionista, el aviso del cambio de metodología de 2023, la
partición de entrenamiento y prueba de los modelos y los créditos de los autores. También aclara
que los datos se usan con fines académicos.
